C01 & A transposição rígida da classificação de risco ex ante impõe barreiras proibitivas a startups e MPMEs & \textbf{contestada} & O texto de 2023 já prevê condições diferenciadas para ME/EPP e startups e sandbox; o efeito sobre a entrada é empírico e não está medido no corpus \\
C02 & O custo de conformidade (relatórios e auditorias) excede o que as empresas locais conseguem absorver & \textbf{sem evidência no corpus} & Sem evidência no corpus. Hipótese a testar: exige estimativa de custo (a análise de impacto do Reino Unido traz um método reaproveitável). \\
C03 & Diante da escassez de profissionais STEM, a burocracia desvia talento técnico para compliance & \textbf{sem evidência no corpus} & Sem evidência no corpus: o PBIA não quantifica escassez STEM. Exige indicador (ContextIndicator) com fonte. \\
C04 & O Brasil não prevê instância consultiva técnica independente equivalente ao painel científico do AI Act & \textbf{contestada} & O texto aprovado pelo Senado em dez/2024 cria um comitê de especialistas e cientistas (CECIA) dentro do SIA; a comparação com o painel do AI Act depende de independência e mandato, não de existência; Válida só para o texto de 2023. Segundo Alanoca et al. (2025), o substitutivo do Senado cria o CECIA: incorporar o texto primário (RF-01). \\
C05 & O Brasil carece de salvaguardas rígidas contra vieses em serviços públicos e no uso estatal de dados & \textbf{contestada} & O PL prevê direito à não discriminação e à correção de vieses e exige, no uso estatal de alto risco, dados testados contra vieses; a lacuna real é que nada disso está em vigor \\
C06 & O marco brasileiro deve evitar proibições rígidas e priorizar sandboxes regulatórios & \textbf{contestada} & Proibição de risco excessivo e sandbox coexistem no PL e no AI Act; escolher entre eles é decisão normativa, não constatação \\
C07 & O PBIA destina cerca de 0,4\% do orçamento ao suporte regulatório e de governança, com risco de estrangulamento institucional & \textbf{com evidência} & Percentual confirmado (0,45\%). O 'risco de estrangulamento' é inferência: depende de quanto a supervisão custaria, o que o corpus não estima. \\
C08 & A matriz elétrica 89,2\% renovável dá vantagem para data centers verdes & \textbf{com evidência} & -- \\
C09 & SUS e CadÚnico são repositórios de dados sem paralelo para treinar modelos & \textbf{com evidência} & Evidência cobre só o SUS. CadÚnico e 'sem paralelo global' não aparecem no corpus. \\
C10 & LLMs em português e em línguas originárias protegem a soberania linguística e cultural & \textbf{com evidência} & O PBIA prevê LLM em português; línguas originárias não aparecem no corpus. \\
C11 & A Nuvem Soberana gerida por Serpro e Dataprev protege dados sensíveis frente a infraestrutura estrangeira & \textbf{com evidência} & -- \\
C12 & A expansão do supercomputador Santos Dumont sustenta a capacidade soberana de computação & \textbf{com evidência} & -- \\
